%% notes-tex/026-kinetics-thermodynamics/sections/2-predictors.tex
%% [[experiments.026-metabolism-flux.kinetic-and-thermodynamic-parameter-datasets]]
%% https://github.com/Mjvolk3/torchcell/tree/main/notes-tex/026-kinetics-thermodynamics/sections/2-predictors.tex

\section{The predictors, and what running them taught\secstatus{todo}}

\subsection{One module per model, on the node-embedding pattern\secstatus{todo}}

Which predictor feeds the flux layer is a design choice to experiment over, not a pipeline
detail to settle once, so the kinetic predictors copy the shape the sequence-embedding
datasets already use: an abstract base, one module per model, and a name-to-class registry
that a config string selects.

One structural difference drives everything else. A node embedding is keyed by
\emph{gene}; a kinetic parameter is keyed by a \emph{pair}, because a turnover number is a
property of one enzyme acting on one substrate. The key is
$(\textrm{uniprot}, \textrm{substrate}, \textrm{predictor}, \textrm{parameter})$.
Collapsing to one row per gene would silently pick a substrate, and collapsing to one row
per catalytic unit would silently pick a subunit, so the store holds pairs and aggregation
lives in the resolver. The dataset records what was predicted; the caller decides what the
layer sees.

\subsection{Sequences come from the genome, not from the model's bundled table\secstatus{todo}}

yeast-GEM ships a \texttt{swissprot.tsv} with a sequence column, and it sits conveniently
beside the model. It is the wrong source. Protein sequences are read from
\texttt{SCerevisiaeGenome} through \texttt{genome[orf].protein.seq}, the same path the ESM2
embedding dataset uses, so everything stays on the derivation chain from the S288C
reference FASTA plus the GAF. The difference is provenance rather than content: a sequence
from a third-party table is a claim by that table's compiler, while the same sequence from
the genome is the one the perturbation ontology edits, so a strain's genotype and the
protein a predictor sees cannot silently disagree.

All 1{,}161 yeast-GEM genes resolve as \texttt{current} through the shared gene-name
reconciler, and all 1{,}161 have a protein sequence, so nothing is lost by the switch.
Rebuilding DLKcat on genome sequences moved the median $k_{cat}$ from 4.0125 to 4.0116
s$^{-1}$, which says the two sources agree on content and differ on provenance.

Two failures are worth recording because both are silent. \texttt{orf in genome} looks
safer than indexing and is not: the genome defines no \texttt{\_\_contains\_\_}, so Python
falls back to integer indexing and calls \texttt{genome[0]}, raising before the real lookup
happens. And \texttt{go\_root} is a separate constructor argument defaulting to a relative
path; left unset it misses and the constructor tries to fetch \texttt{go.obo} from
\texttt{current.geneontology.org}, which answers 403.

\subsection{Validation caught a log base\secstatus{todo}}

The DLKcat implementation reproduces the authors' three released example values to four
decimal places. That check was worth more than its cost, because it caught two defects that
would otherwise have been invisible in a plausible-looking table.

\textbf{The output is $\log_2$, not $\log_{10}$.} The authors exponentiate with
\texttt{math.pow(2, \ldots)}. Reading it as base 10 inflates every turnover number by a
power of roughly 3.3 and still produces numbers that look like credible turnover numbers.

\textbf{Multi-component SMILES are refused by the model itself.} Its fingerprint pass
indexes a node list built only from bonded atoms, so a disconnected fragment, such as the
lone \texttt{[Fe+2]} in the heme SMILES yeast-GEM gives for ferricytochrome $c$, walks off
the end of the list. Excluding them is what the model does, so the exclusion is recorded
rather than worked around.

\subsection{Embedding on four GPUs\secstatus{todo}}

UniKP's protein half is ProtT5-XL, a 3-billion-parameter encoder, and its pinned
environment holds CPU-only torch. The embedding is therefore produced once on GPU and
handed across as a content-addressed cache, keyed by a hash of the sequence so that any
predictor wanting a ProtT5 representation can reuse it and two ORFs with one sequence
resolve to one vector.

Sequences are split by index across the four GPUs and merged; within a shard they are
sorted by length and batched, which is where most of the gain comes from, and a
per-sequence masked mean keeps the result identical to unbatched inference. The 1{,}154
distinct sequences finish in under 40 seconds.

The genome is read exactly once, by a dump step, and the shards read that dump. Four
processes opening the same gffutils SQLite database concurrently raced and made it briefly
unreadable; since the shards need the genome for nothing but a sequence lookup, the
dependency was removed rather than serialized.

\subsection{What each predictor cost to run\secstatus{todo}}

%% SOURCE: hand-authored from the mirror manifests written by
%% experiments/026-metabolism-flux/scripts/mirror_kinetic_models.py; the machine
%% readable form is experiments/026-metabolism-flux/results/kinetic_model_mirrors.json
\begin{table}[H]
  \centering
  \footnotesize
  \caption[Predictor status]{The six predictors named in Wu et al.'s Methods, what each
  emits, and what it took to run here. ``Ships weights'' refers to the repository as
  cloned; UniKP's fitted regressors are released separately on HuggingFace and are
  mirrored alongside it.}
  \label{tab:predictors}
  \begin{tabular}{llll}
    \hline
    Predictor & Emits & Ships weights & Note \\
    \hline
    DLKcat & $k_{cat}$ & yes, extensionless & output is $\log_2$ \\
    UniKP & $k_{cat}$, $K_M$ & regressors on HuggingFace & needs scikit-learn 1.2.2 \\
    TurNuP & $k_{cat}$ & Zenodo archive & takes full reaction SMILES \\
    EITLEM-Kinetics & $k_{cat}$, $K_M$ & Zenodo, 12.3 GB & many checkpoints \\
    Boost\_KM & $K_M$ & yes, two regressors & ESM-1b variant is the one Wu ran \\
    DeepEnzyme & $k_{cat}$ & figshare & uses AlphaFold structures \\
    \hline
  \end{tabular}
\end{table}
